\documentclass[aspectratio=169]{beamer}
\usetheme{metropolis}
\usepackage{booktabs}
\usepackage{listings}
\usepackage{tikz}
\usepackage{hyperref}
\usepackage{xcolor}
\usepackage{graphicx}
\usepackage{fontspec}
\usepackage{emoji}

\usetikzlibrary{shapes.geometric, arrows, positioning}

% Color scheme
\definecolor{darkblue}{RGB}{0,82,155}
\definecolor{lightblue}{RGB}{100,181,246}
\definecolor{darkgreen}{RGB}{46,125,50}
\definecolor{orange}{RGB}{255,152,0}
\definecolor{lightred}{RGB}{239,83,80}

\setbeamercolor{frametitle}{bg=darkblue}
\setbeamercolor{progress bar}{fg=lightblue,bg=lightblue!30}

% Code listing settings
\lstset{
    basicstyle=\ttfamily\small,
    breaklines=true,
    frame=single,
    backgroundcolor=\color{gray!10},
    keywordstyle=\color{darkblue}\bfseries,
    commentstyle=\color{darkgreen}\itshape,
    stringstyle=\color{orange},
    showstringspaces=false,
    language=Python
}

\title{Lecture 4: Data Cleaning \& Preprocessing}
\subtitle{Week 2: Computational Linguistics}
\author{LLM Course}
\date{}

\begin{document}

% Title slide
\begin{frame}
    \titlepage
\end{frame}

% Learning Objectives
\begin{frame}{Learning Objectives 🎯}
    By the end of this lecture, you will be able to:

    \begin{enumerate}
        \item Understand why data cleaning is critical for NLP tasks
        \item Apply common preprocessing techniques to raw text
        \item Use web scraping to collect text data
        \item Implement lemmatization and stemming
        \item Make informed decisions about preprocessing strategies
    \end{enumerate}

    \vspace{1em}
    \textbf{Key theme:} Garbage in, garbage out! 🗑️➡️📊
\end{frame}

% Introduction
\begin{frame}{Why Data Cleaning Matters 📝}
    \textbf{Real-world text is messy!}

    \begin{columns}[T]
        \begin{column}{0.48\textwidth}
            \textbf{Common problems:}
            \begin{itemize}
                \item HTML/XML tags: \texttt{<p>text</p>}
                \item Special characters: \&nbsp; \&amp;
                \item Inconsistent formatting
                \item Extra whitespace
                \item Mixed encodings (UTF-8, ASCII...)
                \item Typos and misspellings
                \item Emoji and Unicode 😊🎉
            \end{itemize}
        \end{column}
        \begin{column}{0.48\textwidth}
            \textbf{Consequences of dirty data:}
            \begin{itemize}
                \item Models learn noise, not signal
                \item Reduced accuracy
                \item Inconsistent predictions
                \item Poor generalization
                \item Wasted computation
            \end{itemize}
        \end{column}
    \end{columns}

    \vspace{1em}
    \begin{center}
        \textit{"Quality of input = quality of output"}
    \end{center}
\end{frame}

\begin{frame}{The Data Cleaning Pipeline 🔄}
    \begin{center}
        \begin{tikzpicture}[
            node distance=1.5cm,
            box/.style={rectangle, draw, fill=lightblue!30, text width=2.5cm, align=center, minimum height=1cm},
            arrow/.style={->, >=stealth, thick}
        ]
            \node[box] (raw) {Raw Text 📄};
            \node[box, right of=raw, xshift=2cm] (clean) {Remove Tags 🏷️};
            \node[box, right of=clean, xshift=2cm] (normalize) {Normalize 🔧};
            \node[box, below of=normalize, yshift=-0.5cm] (tokenize) {Tokenize ✂️};
            \node[box, left of=tokenize, xshift=-2cm] (lemma) {Lemmatize 🌱};
            \node[box, left of=lemma, xshift=-2cm] (final) {Clean Text ✅};

            \draw[arrow] (raw) -- (clean);
            \draw[arrow] (clean) -- (normalize);
            \draw[arrow] (normalize) -- (tokenize);
            \draw[arrow] (tokenize) -- (lemma);
            \draw[arrow] (lemma) -- (final);
        \end{tikzpicture}
    \end{center}

    \vspace{1em}
    \textbf{Note:} Pipeline varies by task! Not all steps are always needed.
\end{frame}

\begin{frame}{Task-Specific Preprocessing 🎯}
    \textbf{Different tasks need different preprocessing:}

    \begin{table}
        \centering
        \small
        \begin{tabular}{lll}
            \toprule
            \textbf{Task} & \textbf{Keep} & \textbf{Remove/Modify} \\
            \midrule
            Sentiment Analysis & Punctuation (!?) & HTML tags, URLs \\
            & Emoji & Extra whitespace \\
            \midrule
            Named Entity & Case (names) & HTML tags \\
            Recognition & Punctuation & Special chars \\
            \midrule
            Text Generation & Everything! & Minimal cleaning \\
            & (preserve style) & (just corrupted text) \\
            \midrule
            Topic Modeling & Content words & Stop words, punctuation \\
            & & Case (lowercase) \\
            \bottomrule
        \end{tabular}
    \end{table}

    \vspace{1em}
    \textbf{Key principle:} Preserve information relevant to your task! 🔑
\end{frame}

% Web Scraping
\begin{frame}{Collecting Data: Web Scraping 🌐}
    \textbf{Why web scraping?}
    \begin{itemize}
        \item Web = massive source of text data
        \item News articles, reviews, social media, forums...
        \item Fresh, diverse, domain-specific content
    \end{itemize}

    \vspace{0.5em}
    \textbf{Popular tools:}
    \begin{itemize}
        \item \textbf{Beautiful Soup:} Parse HTML/XML, extract text
        \item \textbf{Requests:} Fetch web pages
        \item \textbf{Scrapy:} Full-featured web scraping framework
        \item \textbf{Selenium:} For JavaScript-heavy sites
    \end{itemize}

    \vspace{0.5em}
    \textbf{⚠️ Important considerations:}
    \begin{itemize}
        \item Check \texttt{robots.txt} and terms of service
        \item Respect rate limits (don't overwhelm servers)
        \item Be ethical: respect privacy and copyright
    \end{itemize}
\end{frame}

\begin{frame}[fragile]{Web Scraping with Beautiful Soup 🍜}
    \textbf{Basic example:}

    \begin{lstlisting}[language=Python]
from bs4 import BeautifulSoup
import requests

# Fetch webpage
url = "https://example.com/article"
response = requests.get(url)
html_content = response.content

# Parse HTML
soup = BeautifulSoup(html_content, 'html.parser')

# Extract text from specific elements
title = soup.find('h1').get_text()
article = soup.find('article').get_text()

# Remove extra whitespace
clean_text = ' '.join(article.split())

print(f"Title: {title}")
print(f"Article: {clean_text[:200]}...")
    \end{lstlisting}
\end{frame}

\begin{frame}[fragile]{Advanced Beautiful Soup Techniques 🔍}
    \begin{lstlisting}[language=Python]
from bs4 import BeautifulSoup

html = """<div class="article">
    <h2>Breaking News</h2>
    <p class="content">First paragraph.</p>
    <p class="content">Second paragraph.</p>
    <div class="ads">Advertisement</div>
</div>"""

soup = BeautifulSoup(html, 'html.parser')

# Find all paragraphs with class="content"
paragraphs = soup.find_all('p', class_='content')
text = ' '.join([p.get_text() for p in paragraphs])

# Remove unwanted sections
for ad in soup.find_all('div', class_='ads'):
    ad.decompose()  # Delete from tree

# Get clean article text
article_text = soup.get_text(separator=' ', strip=True)
print(article_text)
    \end{lstlisting}
\end{frame}

\begin{frame}{Handling Text Encodings 🔤}
    \textbf{Common encoding issues:}

    \begin{itemize}
        \item UTF-8 vs. ASCII vs. Latin-1
        \item Special characters: é, ñ, 中文
        \item Emoji: 😊 (requires Unicode support)
    \end{itemize}

    \vspace{0.5em}
    \textbf{Best practices:}
    \begin{itemize}
        \item Always use UTF-8 when possible
        \item Detect encoding: \texttt{chardet} library
        \item Normalize Unicode: NFKC vs. NFD forms
    \end{itemize}

    \vspace{0.5em}
    \textbf{Example problem:}
    \begin{center}
        \begin{tabular}{ll}
            \textbf{Input:} & \texttt{caf\textbackslash xe9} (Latin-1) \\
            \textbf{Decoded:} & café (UTF-8) \\
        \end{tabular}
    \end{center}
\end{frame}

% Common Preprocessing Steps
\begin{frame}{Common Preprocessing Steps 🔧}
    \textbf{1. HTML/XML Tag Removal}
    \begin{itemize}
        \item Strip markup: \texttt{<p>Hello</p>} → \texttt{Hello}
        \item Tools: Beautiful Soup, regex, html2text
    \end{itemize}

    \vspace{0.3em}
    \textbf{2. Whitespace Normalization}
    \begin{itemize}
        \item Remove extra spaces, tabs, newlines
        \item \texttt{"Hello\ \ \ world\textbackslash n"} → \texttt{"Hello world"}
    \end{itemize}

    \vspace{0.3em}
    \textbf{3. Punctuation Handling}
    \begin{itemize}
        \item Keep for sentiment: "Great!" vs. "Great"
        \item Remove for topic modeling
    \end{itemize}

    \vspace{0.3em}
    \textbf{4. Case Normalization}
    \begin{itemize}
        \item Lowercase: "Apple" and "apple" → same token
        \item Preserve for NER: "Apple Inc." vs. "apple fruit"
    \end{itemize}

    \vspace{0.3em}
    \textbf{5. Special Characters}
    \begin{itemize}
        \item URLs, email addresses, numbers
        \item Replace or remove based on task
    \end{itemize}
\end{frame}

\begin{frame}[fragile]{Preprocessing Example 💻}
    \begin{lstlisting}[language=Python]
import re

def preprocess_text(text):
    # Remove URLs
    text = re.sub(r'http\S+|www.\S+', '', text)

    # Remove email addresses
    text = re.sub(r'\S+@\S+', '', text)

    # Remove HTML tags
    text = re.sub(r'<.*?>', '', text)

    # Remove extra whitespace
    text = ' '.join(text.split())

    # Lowercase (optional)
    text = text.lower()

    return text

# Example
raw = "Check out https://example.com! <b>Amazing</b>   deals!!"
clean = preprocess_text(raw)
print(clean)  # "check out amazing deals!!"
    \end{lstlisting}
\end{frame}

% Stemming and Lemmatization
\begin{frame}{Reducing Words to Base Forms 🌱}
    \textbf{Goal:} Reduce inflected/derived words to root form

    \begin{columns}[T]
        \begin{column}{0.48\textwidth}
            \textbf{Stemming ✂️}
            \begin{itemize}
                \item Rule-based, crude chopping
                \item Fast, simple
                \item May produce non-words
                \item Porter Stemmer (1980)
                \item Example: "running" → "run"
            \end{itemize}
        \end{column}
        \begin{column}{0.48\textwidth}
            \textbf{Lemmatization 🧠}
            \begin{itemize}
                \item Uses vocabulary + morphology
                \item Slower, more accurate
                \item Produces real words
                \item Requires POS context
                \item Example: "better" → "good"
            \end{itemize}
        \end{column}
    \end{columns}

    \vspace{1em}
    \textbf{When to use which?}
    \begin{itemize}
        \item Stemming: Speed matters, approximate matching OK
        \item Lemmatization: Need interpretable, real words
    \end{itemize}
\end{frame}

\begin{frame}{Stemming vs. Lemmatization Examples ⚖️}
    \begin{table}
        \centering
        \begin{tabular}{lll}
            \toprule
            \textbf{Original Word} & \textbf{Stemming} & \textbf{Lemmatization} \\
            \midrule
            running & run & run \\
            ran & ran & run \\
            runs & run & run \\
            runner & runner & runner \\
            \midrule
            better & better & good \\
            best & best & good \\
            \midrule
            mice & mice & mouse \\
            geese & gees & goose \\
            \midrule
            organization & organ & organization \\
            organizing & organ & organize \\
            \midrule
            meeting & meet & meeting (or meet) \\
            & & (context-dependent!) \\
            \bottomrule
        \end{tabular}
    \end{table}
\end{frame}

\begin{frame}[fragile]{Stemming with NLTK 🔪}
    \begin{lstlisting}[language=Python]
from nltk.stem import PorterStemmer, SnowballStemmer

porter = PorterStemmer()
snowball = SnowballStemmer('english')

words = ['running', 'runs', 'runner', 'ran', 'easily', 'fairly']

print("Word          Porter      Snowball")
print("-" * 40)
for word in words:
    p_stem = porter.stem(word)
    s_stem = snowball.stem(word)
    print(f"{word:12} {p_stem:12} {s_stem}")

# Output:
# Word          Porter      Snowball
# ----------------------------------------
# running      run          run
# runs         run          run
# runner       runner       runner
# ran          ran          ran
# easily       easili       easili
# fairly       fairli       fair
    \end{lstlisting}
\end{frame}

\begin{frame}[fragile]{Lemmatization with spaCy 🚀}
    \begin{lstlisting}[language=Python]
import spacy

nlp = spacy.load("en_core_web_sm")

text = "The cats are running faster than dogs ran yesterday"
doc = nlp(text)

print(f"{'Word':<12} {'Lemma':<12} {'POS'}")
print("-" * 36)
for token in doc:
    print(f"{token.text:<12} {token.lemma_:<12} {token.pos_}")

# Output:
# Word         Lemma        POS
# ------------------------------------
# The          the          DET
# cats         cat          NOUN
# are          be           AUX
# running      run          VERB
# faster       fast         ADV
# than         than         SCONJ
# dogs         dog          NOUN
# ran          run          VERB
# yesterday    yesterday    NOUN
    \end{lstlisting}
\end{frame}

\begin{frame}{spaCy vs. NLTK for Lemmatization 🥊}
    \begin{table}
        \centering
        \begin{tabular}{lll}
            \toprule
            \textbf{Feature} & \textbf{spaCy} & \textbf{NLTK} \\
            \midrule
            Speed & ⚡ Fast (Cython) & 🐢 Slower (Python) \\
            Accuracy & 🎯 High & 👍 Good \\
            Setup & Easy & Requires data download \\
            POS tagging & Built-in & Separate step \\
            Dependencies & Included & Manual WordNet \\
            Use case & Production & Research/Teaching \\
            \bottomrule
        \end{tabular}
    \end{table}

    \vspace{1em}
    \textbf{Recommendation:} Use spaCy for most tasks!

    It's faster, more accurate, and easier to use in production.
\end{frame}

\begin{frame}{Stop Words Removal 🛑}
    \textbf{What are stop words?}
    \begin{itemize}
        \item Common words with little semantic value
        \item Examples: "the", "a", "is", "in", "and", "or"
        \item Language-specific
    \end{itemize}

    \vspace{0.5em}
    \textbf{When to remove?}
    \begin{itemize}
        \item ✅ Topic modeling, text classification
        \item ✅ Search engines, information retrieval
        \item ❌ Sentiment analysis ("not good" vs. "good")
        \item ❌ Machine translation
        \item ❌ Text generation
    \end{itemize}

    \vspace{0.5em}
    \textbf{Warning:} Modern neural models often don't need stop word removal!

    They can learn to ignore them or use them for grammar.
\end{frame}

\begin{frame}[fragile]{Stop Words in Practice 💼}
    \begin{lstlisting}[language=Python]
from nltk.corpus import stopwords
import spacy

# NLTK approach
stop_words_nltk = set(stopwords.words('english'))
text = "This is an example showing stop word removal"
tokens = text.lower().split()
filtered_nltk = [w for w in tokens if w not in stop_words_nltk]

print("NLTK:", filtered_nltk)
# Output: ['example', 'showing', 'stop', 'word', 'removal']

# spaCy approach
nlp = spacy.load("en_core_web_sm")
doc = nlp(text)
filtered_spacy = [token.text for token in doc
                  if not token.is_stop]

print("spaCy:", filtered_spacy)
# Output: ['example', 'showing', 'stop', 'word', 'removal']
    \end{lstlisting}
\end{frame}

% Case Study
\begin{frame}{Case Study: Preprocessing for Sentiment Analysis 📊}
    \textbf{Scenario:} Analyze customer reviews from Amazon

    \vspace{0.5em}
    \textbf{Raw review example:}
    \begin{quote}
        \textit{"I LOVE this product!!! 😍 Best purchase ever! See details at http://example.com"}
    \end{quote}

    \vspace{0.5em}
    \textbf{Preprocessing decisions:}
    \begin{itemize}
        \item ✅ Keep: Punctuation (!), emoji (😍), capitalization (emphasis)
        \item ❌ Remove: URLs, HTML tags
        \item 🤔 Maybe: Extra repeated punctuation (!!!)
    \end{itemize}

    \vspace{0.5em}
    \textbf{Result:}
    \begin{quote}
        \textit{"I LOVE this product! 😍 Best purchase ever!"}
    \end{quote}

    Preserves emotion and intensity!
\end{frame}

\begin{frame}{Discussion: Preprocessing Trade-offs 🤔}
    \textbf{Questions to consider:}

    \begin{enumerate}
        \item \textbf{Information loss:} What do we lose when we lowercase everything?
        \begin{itemize}
            \item Think: "US" (United States) vs. "us" (pronoun)
        \end{itemize}

        \item \textbf{Task dependency:} Why does preprocessing differ by task?
        \begin{itemize}
            \item Hint: What signals matter for sentiment vs. topic modeling?
        \end{itemize}

        \item \textbf{Modern models:} Do transformer models still need heavy preprocessing?
        \begin{itemize}
            \item Consider: BERT handles subwords, capitalization, punctuation...
        \end{itemize}

        \item \textbf{Bias introduction:} Can preprocessing introduce bias?
        \begin{itemize}
            \item Example: Removing slang might remove cultural markers
        \end{itemize}
    \end{enumerate}
\end{frame}

% Practical Tips
\begin{frame}{Practical Tips for Data Cleaning 💡}
    \textbf{Best practices:}

    \begin{enumerate}
        \item \textbf{Inspect your data first!} 👀
        \begin{itemize}
            \item Look at samples before deciding on preprocessing
        \end{itemize}

        \item \textbf{Keep raw data separate} 💾
        \begin{itemize}
            \item Never overwrite originals—you might need them!
        \end{itemize}

        \item \textbf{Document your pipeline} 📝
        \begin{itemize}
            \item Track what preprocessing you applied and why
        \end{itemize}

        \item \textbf{Experiment!} 🧪
        \begin{itemize}
            \item Try different approaches, measure impact on task
        \end{itemize}

        \item \textbf{Validate incrementally} ✅
        \begin{itemize}
            \item Check results after each preprocessing step
        \end{itemize}

        \item \textbf{Consider automation} 🤖
        \begin{itemize}
            \item Use libraries: spaCy, NLTK, HuggingFace tokenizers
        \end{itemize}
    \end{enumerate}
\end{frame}

\begin{frame}[fragile]{Complete Preprocessing Pipeline Example 🔧}
    \begin{lstlisting}[language=Python]
import spacy
import re

class TextPreprocessor:
    def __init__(self, remove_stopwords=False):
        self.nlp = spacy.load("en_core_web_sm")
        self.remove_stopwords = remove_stopwords

    def clean(self, text):
        # Remove URLs
        text = re.sub(r'http\S+', '', text)
        # Remove HTML tags
        text = re.sub(r'<.*?>', '', text)
        # Normalize whitespace
        text = ' '.join(text.split())

        # Process with spaCy
        doc = self.nlp(text)

        # Lemmatize and optionally remove stop words
        tokens = [token.lemma_ for token in doc
                  if not (self.remove_stopwords and token.is_stop)]

        return ' '.join(tokens)
    \end{lstlisting}
\end{frame}

\begin{frame}{Preprocessing Checklist ✅}
    Before starting your project, ask:

    \begin{itemize}
        \item[$\square$] What is my task? (classification, generation, extraction...)
        \item[$\square$] What signals are important? (sentiment, topics, entities...)
        \item[$\square$] Should I lowercase? (preserve case for names?)
        \item[$\square$] How to handle punctuation? (keep for emotion?)
        \item[$\square$] Do I need lemmatization? (or will model handle it?)
        \item[$\square$] Should I remove stop words? (or keep for grammar?)
        \item[$\square$] How to handle special characters? (emoji, numbers, symbols)
        \item[$\square$] What about rare/unknown words? (keep, remove, replace?)
        \item[$\square$] Have I validated on sample data? (inspect before/after!)
    \end{itemize}

    \vspace{0.5em}
    \textbf{Remember:} There's no one-size-fits-all solution! 🎯
\end{frame}

% Resources
\begin{frame}{Tools and Libraries 🛠️}
    \textbf{Web Scraping:}
    \begin{itemize}
        \item Beautiful Soup: \url{https://www.crummy.com/software/BeautifulSoup/}
        \item Scrapy: \url{https://scrapy.org/}
    \end{itemize}

    \vspace{0.3em}
    \textbf{Text Processing:}
    \begin{itemize}
        \item spaCy: \url{https://spacy.io/}
        \item NLTK: \url{https://www.nltk.org/}
        \item TextBlob: \url{https://textblob.readthedocs.io/}
    \end{itemize}

    \vspace{0.3em}
    \textbf{Encoding:}
    \begin{itemize}
        \item chardet: Character encoding detection
        \item ftfy: Fixes broken Unicode
    \end{itemize}

    \vspace{0.3em}
    \textbf{HuggingFace Resources:}
    \begin{itemize}
        \item \href{https://huggingface.co/learn/nlp-course/chapter3/2}{Chapter 3.2: Processing Data}
        \item \href{https://huggingface.co/learn/nlp-course/chapter2/4}{Chapter 2.4: Tokenizers}
    \end{itemize}
\end{frame}

\begin{frame}{Primary References 📚}
    \textbf{Classic papers:}
    \begin{itemize}
        \item Porter, M. F. (1980). An algorithm for suffix stripping. \textit{Program}, 14(3), 130-137.
        \begin{itemize}
            \item The original Porter Stemmer algorithm
        \end{itemize}

        \item Manning, C. D., Raghavan, P., \& Schütze, H. (2008). \textit{Introduction to Information Retrieval}. Cambridge University Press.
        \begin{itemize}
            \item Chapter 2: Text preprocessing fundamentals
        \end{itemize}
    \end{itemize}

    \vspace{0.5em}
    \textbf{Modern resources:}
    \begin{itemize}
        \item spaCy documentation: Industrial-strength NLP
        \item HuggingFace NLP Course: Modern preprocessing with transformers
    \end{itemize}

    \vspace{0.5em}
    \textbf{Ethical considerations:}
    \begin{itemize}
        \item Liang et al. (2020). "Towards Debiasing Sentence Representations"
        \item Consider how preprocessing choices affect fairness
    \end{itemize}
\end{frame}

\begin{frame}{Hands-On Exercise 🧪}
    \textbf{Try this yourself:}

    \begin{enumerate}
        \item Choose a website (news, blog, Reddit...)
        \item Scrape 10-20 articles/posts
        \item Apply different preprocessing pipelines:
        \begin{itemize}
            \item Minimal: Just remove HTML
            \item Moderate: + normalize whitespace, lowercase
            \item Heavy: + lemmatize, remove stop words
        \end{itemize}
        \item Compare the results:
        \begin{itemize}
            \item How does vocabulary size change?
            \item What information is lost/preserved?
            \item Which would work best for sentiment analysis? Topic modeling?
        \end{itemize}
    \end{enumerate}

    \vspace{0.5em}
    \textbf{Bonus:} Share interesting findings with classmates! 🎉
\end{frame}

% Summary
\begin{frame}{Key Takeaways 🎯}
    \begin{enumerate}
        \item \textbf{Preprocessing is crucial} but task-dependent
        \begin{itemize}
            \item No universal pipeline—adapt to your needs!
        \end{itemize}

        \item \textbf{Balance cleaning vs. information loss}
        \begin{itemize}
            \item More preprocessing ≠ always better
        \end{itemize}

        \item \textbf{Stemming vs. Lemmatization}
        \begin{itemize}
            \item Stemming: fast, approximate
            \item Lemmatization: slow, accurate
        \end{itemize}

        \item \textbf{Modern models are robust}
        \begin{itemize}
            \item Transformers can handle messy text better than older models
        \end{itemize}

        \item \textbf{Always validate}
        \begin{itemize}
            \item Inspect data before and after preprocessing
        \end{itemize}
    \end{enumerate}

    \vspace{1em}
    \textbf{Next lecture:} Tokenization deep dive! 🔤
\end{frame}

\begin{frame}[standout]
    Questions? 🤔

    \vspace{2em}

    \normalsize
    Next: Lecture 5 - Tokenization
\end{frame}

\end{document}
